Selective homomorphic encryption aims to reduce the high cost of fully homomorphic encryption by encrypting only the sensitive part of the data. Region-of-interest (ROI) based approaches proposed in 2026 encrypt only the ROI of a medical image and process the rest in plaintext. They prove their security with respect to a leakage function that reveals the unencrypted region and the location, size and shape of the ROI to the server. This thesis investigates how much this permitted leakage reveals about the clinical content of the encrypted region in real medical images, and what it costs to remove the leakage.

The $\Pi_{\mathrm{ROI}}$ protocol was reproduced with TenSEAL CKKS using the same parameters. Two leakage-abuse attacks were developed on chest X-ray (COVID-QU-Ex, 33,920 images) and brain MRI (3,064 slices, 233 patients) data.
\begin{itemize}
\item An attack that uses only the ROI metadata predicts the tumor type with a macro AUC of $0.83$ and pneumonia with an AUC of $0.90$.
\item A convolutional neural network attacker that sees the image with the ROI hidden recovers the diagnosis with an AUC of $0.97$--$0.99$. It still reaches an AUC of $0.96$ when 87\% of the image is encrypted.
\item An attacker trained on data from a different source achieves similar performance.
\item Simple preprocessing steps do not remove the leakage.
\end{itemize}

Against an adaptive attacker, reducing the attack AUC below $0.8$ required encrypting 98\% of the chest X-ray and the entire brain MRI. Under this condition, the speedup on $512 \times 512$ images drops from 4.2 times to 1.02 times for chest X-rays and from 24 times to 1.00 times for brain MRI, compared with the default ROI.

The findings show that ROI-based selective homomorphic encryption loses its efficiency advantage on medical images once a privacy requirement is imposed. Based on these results, the thesis proposes design rules that rely on measuring the leakage.

\textbf{Keywords:} Homomorphic Encryption, CKKS, Selective Encryption, Leakage-Abuse Attacks, Medical Image Privacy
